% 期刊论文骨架（英文）。用 xelatex 编译：
% python scripts/compile_latex.py main.tex
%
% 默认用通用 article 类。投稿前按目标期刊换成官方模板：
% - IEEE 系期刊：\documentclass[journal]{IEEEtran}，配合 IEEEtran.bst
% - Elsevier 期刊：\documentclass[review]{elsarticle}
% - Springer 期刊：\documentclass[sn-basic]{sn-jnl}
% 这些类多数 TeX Live 发行版自带。官方模板以期刊 Author Guide 为准，不要手改。
\documentclass[11pt]{article}

\usepackage[margin=2.5cm]{geometry}
\usepackage{amsmath,amssymb}
\usepackage{graphicx}
\usepackage{booktabs}
\usepackage{natbib}
\usepackage[hidelinks]{hyperref}

\title{期刊论文标题}
\author{作者姓名 \\ 机构，城市，国家 \\ \texttt{email@example.com}}
\date{}

\begin{document}

\maketitle

\begin{abstract}
摘要四句：问题、方法、结果、结论。
\end{abstract}

\noindent\textbf{Keywords:} keyword1, keyword2, keyword3

\section{Introduction}
\label{sec:intro}
背景与问题。相关工作线索。本文贡献。

\section{Related Work}
\label{sec:related}
按主题组织，每个主题给出关键文献与差异。

\section{Method}
\label{sec:method}
公式与算法描述。

\section{Experiments}
\label{sec:experiments}
数据、设置、结果。

\begin{table}[htbp]
\centering
\begin{tabular}{lcc}
\toprule
方法 & 指标A & 指标B \\
\midrule
基线 & 0.00 & 0.00 \\
本文 & 0.00 & 0.00 \\
\bottomrule
\end{tabular}
\caption{结果对比。数字须与数据源一致。}
\label{tab:results}
\end{table}

\section{Conclusion}
\label{sec:conclusion}
总结与展望。

\section*{Acknowledgements}
致谢。

% 正式投稿时把 thebibliography 换成 .bib 文件：
% \bibliographystyle{plainnat}
% \bibliography{references}
\begin{thebibliography}{9}
\bibitem[Author et al.(2026)]{example2026} A.~Author, B.~Author. Title of the paper. \textit{Journal Name}, 2026.
\end{thebibliography}

\appendix
\section{Appendix}
\label{app:detail}
补充细节。

\end{document}
